\documentclass[12pt]{article}
\usepackage{fontspec}
\usepackage{fullpage}
\usepackage{hyperref}
\hypersetup{bookmarks=true,colorlinks=true,linkcolor=red,citecolor=blue,filecolor=magenta,urlcolor=cyan}
\usepackage{amsmath}
\usepackage{amssymb}
\usepackage{mathtools}
\usepackage{unicode-math}
\usepackage{tabularray}
\usepackage{tabularx}
\usepackage{booktabs}
\usepackage{caption}
\usepackage{graphics}
\usepackage{svg}
\usepackage{float}
\usepackage{enumitem}
\usepackage{adjustbox}
\usepackage{filecontents}
\usepackage[backend=bibtex]{biblatex}
\usepackage{url}
\setmathfont{Latin Modern Math}
\newcommand{\gt}{\ensuremath >}
\newcommand{\lt}{\ensuremath <}
\newlist{symbDescription}{description}{1}
\setlist[symbDescription]{noitemsep, topsep=0pt, parsep=0pt, partopsep=0pt}
\newcommand{\resizeExpression}[1]{
  \adjustbox{max width=\linewidth}{$#1$}
}
\bibliography{bibfile}
\title{Software Requirements Specification for Charged Particle Trajectory Simulator}
\author{Zhuo Zhang}
\begin{document}
\maketitle
\tableofcontents
\newpage
\section{Reference Material}
\label{Sec:RefMat}
This section records information for easy reference.

\subsection{Table of Units}
\label{Sec:ToU}
The unit system used throughout is SI (Système International d'Unités). In addition to the basic units, several derived units are also used. For each unit, the \hyperref[Table:ToU]{Table of Units} lists the symbol, a description, and the SI name.

\begin{longtblr}
[caption={Table of Units}]
{colspec={l l l}, rowhead=1, hline{1,Z}=\heavyrulewidth, hline{2}=\lightrulewidth}
\textbf{Symbol} & \textbf{Description} & \textbf{SI Name}
\\
${\text{C}}$ & electric charge & coulomb
\\
${\text{kg}}$ & mass & kilogram
\\
${\text{m}}$ & length & metre
\\
${\text{N}}$ & force & newton
\\
${\text{s}}$ & time & second
\\
${\text{T}}$ & magnetic flux density & tesla
\label{Table:ToU}
\end{longtblr}
\subsection{Table of Symbols}
\label{Sec:ToS}
The symbols used in this document are summarized in the \hyperref[Table:ToS]{Table of Symbols} along with their units. Throughout the document, symbols in bold will represent vectors, and scalars otherwise. The symbols are listed in alphabetical order. For vector quantities, the units shown are for each component of the vector.

\begin{longtblr}
[caption={Table of Symbols}]
{colspec={l l l}, rowhead=1, hline{1,Z}=\heavyrulewidth, hline{2}=\lightrulewidth}
\textbf{Symbol} & \textbf{Description} & \textbf{Units}
\\
${a_{\text{x}}}$ & X-acceleration of the particle & $\frac{\text{m}}{\text{s}^{2}}$
\\
${a_{\text{y}}}$ & Y-acceleration of the particle & $\frac{\text{m}}{\text{s}^{2}}$
\\
$\symbf{a}\text{(}t\text{)}$ & Acceleration & $\frac{\text{m}}{\text{s}^{2}}$
\\
${B_{0}}$ & Magnetic flux density in region 0 & ${\text{T}}$
\\
${B_{1}}$ & Magnetic flux density in region 1 & ${\text{T}}$
\\
${B_{2}}$ & Magnetic flux density in region 2 & ${\text{T}}$
\\
${B_{3}}$ & Magnetic flux density in region 3 & ${\text{T}}$
\\
${B_{4}}$ & Magnetic flux density in region 4 & ${\text{T}}$
\\
${B_{5}}$ & Magnetic flux density in region 5 & ${\text{T}}$
\\
$\symbf{B}$ & Magnetic flux density vector & ${\text{T}}$
\\
${\symbf{B}_{i}}$ & Magnetic flux density in region i & ${\text{T}}$
\\
${d_{\text{orient}}}$ & Detector orientation & --
\\
${E_{\text{x}0}}$ & X-electric field in region 0 & $\frac{\text{N}}{\text{C}}$
\\
${E_{\text{x}1}}$ & X-electric field in region 1 & $\frac{\text{N}}{\text{C}}$
\\
${E_{\text{x}2}}$ & X-electric field in region 2 & $\frac{\text{N}}{\text{C}}$
\\
${E_{\text{x}3}}$ & X-electric field in region 3 & $\frac{\text{N}}{\text{C}}$
\\
${E_{\text{x}4}}$ & X-electric field in region 4 & $\frac{\text{N}}{\text{C}}$
\\
${E_{\text{x}5}}$ & X-electric field in region 5 & $\frac{\text{N}}{\text{C}}$
\\
${E_{\text{y}0}}$ & Y-electric field in region 0 & $\frac{\text{N}}{\text{C}}$
\\
${E_{\text{y}1}}$ & Y-electric field in region 1 & $\frac{\text{N}}{\text{C}}$
\\
${E_{\text{y}2}}$ & Y-electric field in region 2 & $\frac{\text{N}}{\text{C}}$
\\
${E_{\text{y}3}}$ & Y-electric field in region 3 & $\frac{\text{N}}{\text{C}}$
\\
${E_{\text{y}4}}$ & Y-electric field in region 4 & $\frac{\text{N}}{\text{C}}$
\\
${E_{\text{y}5}}$ & Y-electric field in region 5 & $\frac{\text{N}}{\text{C}}$
\\
$\symbf{E}$ & Electric field vector & $\frac{\text{N}}{\text{C}}$
\\
${\symbf{E}_{i}}$ & Electric field in region i & $\frac{\text{N}}{\text{C}}$
\\
$\symbf{F}$ & Force & ${\text{N}}$
\\
$h$ & Region height & ${\text{m}}$
\\
${L_{\text{det}}}$ & Detector line & --
\\
$m$ & Particle mass & ${\text{kg}}$
\\
${m_{\text{max}}}$ & Maximum particle mass & ${\text{kg}}$
\\
${m_{\text{min}}}$ & Minimum particle mass & ${\text{kg}}$
\\
$N$ & Number of field regions & --
\\
${N_{\text{col}}}$ & Number of grid columns & --
\\
$\symbf{p}\text{(}t\text{)}$ & Position & ${\text{m}}$
\\
$q$ & Particle charge & ${\text{C}}$
\\
${q_{\text{max}}}$ & Maximum charge magnitude & ${\text{C}}$
\\
${R_{i}}$ & Rectangular field region & --
\\
$s$ & Dependent variables & ${\text{m}}$
\\
${s_{0}}$ & Initial state vector & --
\\
$t$ & Time & ${\text{s}}$
\\
${t_{\text{final}}}$ & Final simulation time & ${\text{s}}$
\\
${t_{\text{hit}}}$ & Time of detector hit & ${\text{s}}$
\\
${t_{\text{max}}}$ & Maximum simulation time & ${\text{s}}$
\\
${v_{\text{max}}}$ & Maximum initial speed & $\frac{\text{m}}{\text{s}}$
\\
${v_{\text{x}}}$ & X-velocity of the particle & $\frac{\text{m}}{\text{s}}$
\\
${v_{\text{x}0}}$ & Initial x-velocity & $\frac{\text{m}}{\text{s}}$
\\
${v_{\text{y}}}$ & Y-velocity of the particle & $\frac{\text{m}}{\text{s}}$
\\
${v_{\text{y}0}}$ & Initial y-velocity & $\frac{\text{m}}{\text{s}}$
\\
$\symbf{v}\text{(}t\text{)}$ & Velocity & $\frac{\text{m}}{\text{s}}$
\\
$\symbf{v}\text{×}\symbf{B}$ & Cross product of velocity and magnetic field & --
\\
$w$ & Region width & ${\text{m}}$
\\
$x$ & X-position of the particle & ${\text{m}}$
\\
${x_{0}}$ & Initial x-position & ${\text{m}}$
\\
${x_{\text{det}}}$ & Detector line x-position & ${\text{m}}$
\\
${x_{\text{detMax}}}$ & Maximum x-coordinate of detector & ${\text{m}}$
\\
${x_{\text{detMin}}}$ & Minimum x-coordinate of detector & ${\text{m}}$
\\
${x_{\text{grid}}}$ & Grid origin x-coordinate & ${\text{m}}$
\\
${x_{\text{hit}}}$ & X-coordinate of impact point & ${\text{m}}$
\\
$y$ & Y-position of the particle & ${\text{m}}$
\\
${y_{0}}$ & Initial y-position & ${\text{m}}$
\\
${y_{\text{det}}}$ & Detector line y-position & ${\text{m}}$
\\
${y_{\text{detMax}}}$ & Maximum y-coordinate of detector & ${\text{m}}$
\\
${y_{\text{detMin}}}$ & Minimum y-coordinate of detector & ${\text{m}}$
\\
${y_{\text{grid}}}$ & Grid origin y-coordinate & ${\text{m}}$
\\
${y_{\text{hit}}}$ & Y-coordinate of impact point & ${\text{m}}$
\\
$κ$ & Charge-to-mass ratio & $\frac{\text{C}}{\text{kg}}$
\label{Table:ToS}
\end{longtblr}
\subsection{Abbreviations and Acronyms}
\label{Sec:TAbbAcc}
\begin{longtblr}
[caption={Abbreviations and Acronyms}]
{colspec={l l}, rowhead=1, hline{1,Z}=\heavyrulewidth, hline{2}=\lightrulewidth}
\textbf{Abbreviation} & \textbf{Full Form}
\\
A & Assumption
\\
DD & Data Definition
\\
GD & General Definition
\\
GS & Goal Statement
\\
IM & Instance Model
\\
LC & Likely Change
\\
PS & Physical System Description
\\
R & Requirement
\\
RefBy & Referenced by
\\
Refname & Reference Name
\\
SRS & Software Requirements Specification
\\
TM & Theoretical Model
\\
UC & Unlikely Change
\\
Uncert. & Typical Uncertainty
\label{Table:TAbbAcc}
\end{longtblr}
\section{Introduction}
\label{Sec:Intro}
The motion of a charged particle through electromagnetic fields is a fundamental phenomenon arising in particle accelerators, mass spectrometers, and plasma confinement devices. It is therefore useful to have a program to simulate charged particle trajectories in piecewise-uniform fields.

The following section provides an overview of the software requirements specification (SRS) for Charged Particle Trajectory Simulator. This section explains the purpose of this document, the scope of the requirements, the characteristics of the intended reader, and the organization of the document.

\subsection{Purpose of Document}
\label{Sec:DocPurpose}
The primary purpose of this document is to record the requirements of Trajecto. Goals, assumptions, theoretical models, definitions, and other model derivation information are specified, allowing the reader to fully understand and verify the purpose and scientific basis of Trajecto. With the exception of \hyperref[Sec:SysConstraints]{system constraints}, this SRS will remain abstract, describing what problem is being solved, but not how to solve it.

This document will be used as a starting point for subsequent development phases, including writing the design specification and the software verification and validation plan. The design document will show how the requirements are to be realized, including decisions on the numerical algorithms and programming environment. The verification and validation plan will show the steps that will be used to increase confidence in the software documentation and the implementation. Although the SRS fits in a series of documents that follow the so-called waterfall model, the actual development process is not constrained in any way. Even when the waterfall model is not followed, as Parnas and Clements point out \cite{parnasClements1986}, the most logical way to present the documentation is still to ``fake'' a rational design process.

\subsection{Scope of Requirements}
\label{Sec:ReqsScope}
The scope of the requirements includes the analysis of two-dimensional charged particle motion in the x-y plane under the Lorentz force.

\subsection{Characteristics of Intended Reader}
\label{Sec:ReaderChars}
Reviewers of this documentation should have an understanding of undergraduate level physics and undergraduate level calculus. The users of Trajecto can have a lower level of expertise, as explained in \hyperref[Sec:UserChars]{Sec:User Characteristics}.

\subsection{Organization of Document}
\label{Sec:DocOrg}
The organization of this document follows the template for an SRS for scientific computing software proposed by \cite{koothoor2013}, \cite{smithLai2005}, \cite{smithEtAl2007}, and \cite{smithKoothoor2016}. The presentation follows the standard pattern of presenting goals, theories, definitions, and assumptions. For readers that would like a more bottom up approach, they can start reading the \hyperref[Sec:IMs]{instance models} and trace back to find any additional information they require.

The \hyperref[Sec:GoalStmt]{goal statements} are refined to the theoretical models and the \hyperref[Sec:TMs]{theoretical models} to the \hyperref[Sec:IMs]{instance models}.

\section{General System Description}
\label{Sec:GenSysDesc}
This section provides general information about the system. It identifies the interfaces between the system and its environment, describes the user characteristics, and lists the system constraints.

\subsection{System Context}
\label{Sec:SysContext}
\hyperref[Figure:sysCtxDiag]{Fig:sysCtxDiag} shows the system context. A rectangle represents the software system (Trajecto), and a circle represents the user providing inputs and receiving outputs.

\begin{figure}[H]
\begin{center}
\includegraphics[width=\textwidth]{../../../../datafiles/trajecto/SystemContextFigure.png}
\caption{System Context}
\label{Figure:sysCtxDiag}
\end{center}
\end{figure}
The user provides the particle properties, field region grid, per-region electromagnetic fields, detector specification, and simulation time; the software system computes and reports the particle trajectory and detector hit outcome:

\subsection{User Characteristics}
\label{Sec:UserChars}
The end user of Trajecto should have an understanding of undergraduate-level physics.

\subsection{System Constraints}
\label{Sec:SysConstraints}
There are no system constraints.

\section{Specific System Description}
\label{Sec:SpecSystDesc}
This section first presents the problem description, which gives a high-level view of the problem to be solved. This is followed by the solution characteristics specification, which presents the assumptions, theories, and definitions that are used.

\subsection{Problem Description}
\label{Sec:ProbDesc}
A system is needed to predict the trajectory of a charged particle moving through piecewise-uniform electric field and magnetic field regions.

\subsubsection{Terminology and Definitions}
\label{Sec:TermDefs}
This subsection provides a list of terms that are used in the subsequent sections and their meaning, with the purpose of reducing ambiguity and making it easier to correctly understand the requirements.

\begin{itemize}
\item{Charged particle: A particle that has a non-zero electric charge and can interact with electric and magnetic fields.}
\item{Electric field: A physical field that exerts an electric force on charged particles.}
\item{Magnetic field: A physical field that exerts a magnetic force on moving charged particles.}
\item{Lorentz force: The force on a charged particle due to electric and magnetic fields.}
\item{Trajectory: The path traced by a particle over time.}
\item{Field region: A spatial region in which the electric field and magnetic field are specified, typically treated as uniform within the region.}
\item{Region grid: A rectangular arrangement of N equal-sized rectangular field regions tiled side-by-side with no gaps or overlaps, whose union forms a single axis-aligned rectangle.}
\item{Detector line: A line segment (either horizontal or vertical) at a specified location where the particle impact position is recorded.}
\item{Velocity selector: A configuration of crossed electric and magnetic fields that allows only particles with a particular velocity to pass through without deflection.}
\item{Cartesian coordinate system: A coordinate system that specifies each point uniquely in a plane by a set of numerical coordinates, which are the signed distances to the point from two fixed perpendicular oriented lines, measured in the same unit of length.}
\end{itemize}
\subsubsection{Physical System Description}
\label{Sec:PhysSyst}
The physical system of Trajecto, as shown in \hyperref[Figure:trajectoPhysSys]{Fig:trajectoPhysSys}, includes the following elements:

\begin{description}[font=\normalfont]
\item[PS1:]{A charged particle (the simulated particle)}
\item[PS2:]{The electric field in each field region (may differ between regions)}
\item[PS3:]{The magnetic field perpendicular to the x-y plane in each field region}
\item[PS4:]{The region grid of N equal-sized rectangular field regions}
\item[PS5:]{The detector line (horizontal or vertical) at a specified location}
\end{description}
\begin{figure}[H]
\begin{center}
\includegraphics[width=0.6\textwidth]{../../../../datafiles/trajecto/trajecto.png}
\caption{The physical system description}
\label{Figure:trajectoPhysSys}
\end{center}
\end{figure}
\subsubsection{Goal Statements}
\label{Sec:GoalStmt}
Given the particle properties ($q$ and $m$), initial conditions (${x_{0}}$, ${y_{0}}$, ${v_{\text{x}0}}$, ${v_{\text{y}0}}$), the number of field regions $N$, the per-region fields (${E_{\text{x}0}}$, ${E_{\text{y}0}}$, ${B_{0}}$ for each region), the region grid specification (${x_{\text{grid}}}$, ${y_{\text{grid}}}$, $w$, $h$), the detector specification (${d_{\text{orient}}}$, ${x_{\text{det}}}$, ${y_{\text{det}}}$ and extent), and the simulation time ${t_{\text{final}}}$, the goal statements are:

\begin{description}[font=\normalfont]
\item[predictTrajectory:\phantomsection\label{predictTrajectory}]{Predict the trajectory of a charged particle under prescribed electric and magnetic fields, computed using \hyperref[IM:stateEvol]{IM:stateEvol}.}
\item[determineDetectorOutcome:\phantomsection\label{determineDetectorOutcome}]{Determine impact position and time-of-flight at a specified detector line, computed using \hyperref[IM:detHit]{IM:detHit}.}
\end{description}
\subsection{Solution Characteristics Specification}
\label{Sec:SolCharSpec}
The instance models that govern Trajecto are presented in the \hyperref[Sec:IMs]{Instance Model Section}. The information to understand the meaning of the instance models and their derivation is also presented, so that the instance models can be verified.

\subsubsection{Assumptions}
\label{Sec:Assumps}
This section simplifies the original problem and helps in developing the theoretical models by filling in the missing information for the physical system. The assumptions refine the scope by providing more detail.

\begin{description}[font=\normalfont]
\item[piecewiseUniform:\phantomsection\label{piecewiseUniform}]{The electric and magnetic fields are uniform within each field region and may change only at region boundaries. (RefBy: \hyperref[DD:fieldsByRegion]{DD:fieldsByRegion}.)}
\item[singleParticle:\phantomsection\label{singleParticle}]{The particle is treated as a point mass and point charge. (RefBy: \hyperref[IM:stateEvol]{IM:stateEvol}, \hyperref[TM:lorentzForce]{TM:lorentzForce}, \hyperref[TM:eqnMotion]{TM:eqnMotion}, and \hyperref[noInteractions]{A:noInteractions}.)}
\item[noInteractions:\phantomsection\label{noInteractions}]{Collisions and particle-particle interactions (including space-charge effects) are neglected; this follows from \hyperref[singleParticle]{A:singleParticle}. (RefBy: \hyperref[IM:stateEvol]{IM:stateEvol}.)}
\item[prescribedFields:\phantomsection\label{prescribedFields}]{The electric and magnetic fields are user-specified and remain fixed during the simulation. (RefBy: \hyperref[TM:lorentzForce]{TM:lorentzForce}, \hyperref[TM:eqnMotion]{TM:eqnMotion}, \hyperref[GD:dyn2D]{GD:dyn2D}, \hyperref[DD:eField]{DD:eField}, \hyperref[DD:fieldsByRegion]{DD:fieldsByRegion}, and \hyperref[DD:bField]{DD:bField}.)}
\item[twoDMotion:\phantomsection\label{twoDMotion}]{The particle motion is confined to the x-y plane. (RefBy: \hyperref[GD:vCrossB2D]{GD:vCrossB2D}, \hyperref[IM:stateEvol]{IM:stateEvol}, \hyperref[GD:kin2D]{GD:kin2D}, \hyperref[GD:dyn2D]{GD:dyn2D}, \hyperref[lcExtendTo3D]{LC:lcExtendTo3D}, \hyperref[eAxisAligned]{A:eAxisAligned}, and \hyperref[bPerpPlane]{A:bPerpPlane}.)}
\item[bPerpPlane:\phantomsection\label{bPerpPlane}]{The magnetic field is perpendicular to the x-y plane, consistent with \hyperref[twoDMotion]{A:twoDMotion}. (RefBy: \hyperref[GD:vCrossB2D]{GD:vCrossB2D}, \hyperref[GD:dyn2D]{GD:dyn2D}, and \hyperref[DD:bField]{DD:bField}.)}
\item[eAxisAligned:\phantomsection\label{eAxisAligned}]{The electric field lies in the x-y plane and is aligned with a coordinate axis, consistent with \hyperref[twoDMotion]{A:twoDMotion}. (RefBy: \hyperref[GD:dyn2D]{GD:dyn2D} and \hyperref[DD:eField]{DD:eField}.)}
\item[rectRegions:\phantomsection\label{rectRegions}]{All field regions are axis-aligned rectangles of identical width w and height h. The N regions are tiled adjacently (sharing edges with no gaps or overlaps) so that their union forms a single axis-aligned rectangle. (RefBy: \hyperref[gridLayout]{A:gridLayout} and \hyperref[DD:regionRect]{DD:regionRect}.)}
\item[gridLayout:\phantomsection\label{gridLayout}]{The field regions are numbered 1 through N and arranged in a single row (left-to-right) or single column (bottom-to-top) within the grid rectangle. The arrangement direction is determined by the region grid specification, and region geometry is constrained by \hyperref[rectRegions]{A:rectRegions}. (RefBy: \hyperref[DD:regionRect]{DD:regionRect}.)}
\item[lineDetector:\phantomsection\label{lineDetector}]{The detector is modelled as a line segment that is either horizontal or vertical, located at the boundary of or within the region grid. (RefBy: \hyperref[IM:detHit]{IM:detHit} and \hyperref[DD:detectorLine]{DD:detectorLine}.)}
\item[fullDetection:\phantomsection\label{fullDetection}]{The detector line is sufficiently long to record the impact point for any trajectory within the scope of the simulation. (RefBy: \hyperref[IM:detHit]{IM:detHit}.)}
\item[lorentzOnly:\phantomsection\label{lorentzOnly}]{The particle dynamics are governed only by the Lorentz force; all other forces are neglected. (RefBy: \hyperref[IM:stateEvol]{IM:stateEvol}, \hyperref[TM:eqnMotion]{TM:eqnMotion}, and \hyperref[ucLorentzForce]{UC:ucLorentzForce}.)}
\end{description}
\subsubsection{Theoretical Models}
\label{Sec:TMs}
This section focuses on the general equations and laws that Trajecto is based on.

\medskip
\noindent
\begin{minipage}{\textwidth}
\begin{tabular}{>{\raggedright}p{0.13\textwidth}>{\raggedright\arraybackslash}p{0.82\textwidth}}
\toprule \textbf{Refname} & \textbf{TM:velocity}
\phantomsection 
\label{TM:velocity}
\\ \midrule
Label & Velocity
        
\\ \midrule
Equation & \begin{displaymath}
           \resizeExpression{\symbf{v}\text{(}t\text{)}=\frac{\,d\symbf{p}\text{(}t\text{)}}{\,dt}}
           \end{displaymath}
\\ \midrule
Description & \begin{symbDescription}
              \item{$\symbf{v}\text{(}t\text{)}$ is the velocity ($\frac{\text{m}}{\text{s}}$)}
              \item{$t$ is the time (${\text{s}}$)}
              \item{$\symbf{p}\text{(}t\text{)}$ is the position (${\text{m}}$)}
              \end{symbDescription}
\\ \midrule
Source & \cite{velocityWiki}
         
\\ \midrule
RefBy & \hyperref[GD:kin2D]{GD:kin2D}
        
\\ \bottomrule
\end{tabular}
\end{minipage}

\medskip
\noindent
\begin{minipage}{\textwidth}
\begin{tabular}{>{\raggedright}p{0.13\textwidth}>{\raggedright\arraybackslash}p{0.82\textwidth}}
\toprule \textbf{Refname} & \textbf{TM:acceleration}
\phantomsection 
\label{TM:acceleration}
\\ \midrule
Label & Acceleration
        
\\ \midrule
Equation & \begin{displaymath}
           \resizeExpression{\symbf{a}\text{(}t\text{)}=\frac{\,d\symbf{v}\text{(}t\text{)}}{\,dt}}
           \end{displaymath}
\\ \midrule
Description & \begin{symbDescription}
              \item{$\symbf{a}\text{(}t\text{)}$ is the acceleration ($\frac{\text{m}}{\text{s}^{2}}$)}
              \item{$t$ is the time (${\text{s}}$)}
              \item{$\symbf{v}\text{(}t\text{)}$ is the velocity ($\frac{\text{m}}{\text{s}}$)}
              \end{symbDescription}
\\ \midrule
Source & \cite{accelerationWiki}
         
\\ \midrule
RefBy & 
\\ \bottomrule
\end{tabular}
\end{minipage}

\medskip
\noindent
\begin{minipage}{\textwidth}
\begin{tabular}{>{\raggedright}p{0.13\textwidth}>{\raggedright\arraybackslash}p{0.82\textwidth}}
\toprule \textbf{Refname} & \textbf{TM:lorentzForce}
\phantomsection 
\label{TM:lorentzForce}
\\ \midrule
Label & Lorentz force law
        
\\ \midrule
Equation & \begin{displaymath}
           \resizeExpression{\symbf{F}=q\,\left(\symbf{E}+\symbf{v}\text{×}\symbf{B}\right)}
           \end{displaymath}
\\ \midrule
Description & \begin{symbDescription}
              \item{$\symbf{F}$ is the force (${\text{N}}$)}
              \item{$q$ is the particle charge (${\text{C}}$)}
              \item{$\symbf{E}$ is the electric field vector ($\frac{\text{N}}{\text{C}}$)}
              \item{$\symbf{v}\text{×}\symbf{B}$ is the cross product of velocity and magnetic field (Unitless)}
              \end{symbDescription}
\\ \midrule
Notes & The electromagnetic force on a charged particle is F = q*(E + v×B), where $\symbf{F}$ is the total force vector, $q$ is the particle charge, $\symbf{E}$ is the electric field vector, and $\symbf{v}\text{×}\symbf{B}$ is the cross product of velocity and magnetic flux density.. For 2D planar motion with an out-of-plane magnetic field, this reduces to the component equations derived in the general definitions. Assumes \hyperref[singleParticle]{A:singleParticle}, \hyperref[prescribedFields]{A:prescribedFields}.
        
\\ \midrule
Source & --
         
\\ \midrule
RefBy & \hyperref[TM:eqnMotion]{TM:eqnMotion} and \hyperref[GD:dyn2D]{GD:dyn2D}
        
\\ \bottomrule
\end{tabular}
\end{minipage}

\medskip
\noindent
\begin{minipage}{\textwidth}
\begin{tabular}{>{\raggedright}p{0.13\textwidth}>{\raggedright\arraybackslash}p{0.82\textwidth}}
\toprule \textbf{Refname} & \textbf{TM:eqnMotion}
\phantomsection 
\label{TM:eqnMotion}
\\ \midrule
Label & Equation of motion under electromagnetic fields
        
\\ \midrule
Equation & \begin{displaymath}
           \resizeExpression{{a_{\text{x}}}=\frac{q}{m}\,\left({E_{\text{x}0}}+{v_{\text{y}}}\,{B_{0}}\right)}
           \end{displaymath}
           \begin{displaymath}
           \resizeExpression{{a_{\text{y}}}=\frac{q}{m}\,\left({E_{\text{y}0}}-{v_{\text{x}}}\,{B_{0}}\right)}
           \end{displaymath}
\\ \midrule
Description & \begin{symbDescription}
              \item{${a_{\text{x}}}$ is the x-acceleration of the particle ($\frac{\text{m}}{\text{s}^{2}}$)}
              \item{$q$ is the particle charge (${\text{C}}$)}
              \item{$m$ is the particle mass (${\text{kg}}$)}
              \item{${E_{\text{x}0}}$ is the x-electric field in region 0 ($\frac{\text{N}}{\text{C}}$)}
              \item{${v_{\text{y}}}$ is the y-velocity of the particle ($\frac{\text{m}}{\text{s}}$)}
              \item{${B_{0}}$ is the magnetic flux density in region 0 (${\text{T}}$)}
              \end{symbDescription}
              \begin{symbDescription}
              \item{${a_{\text{y}}}$ is the y-acceleration of the particle ($\frac{\text{m}}{\text{s}^{2}}$)}
              \item{$q$ is the particle charge (${\text{C}}$)}
              \item{$m$ is the particle mass (${\text{kg}}$)}
              \item{${E_{\text{y}0}}$ is the y-electric field in region 0 ($\frac{\text{N}}{\text{C}}$)}
              \item{${v_{\text{x}}}$ is the x-velocity of the particle ($\frac{\text{m}}{\text{s}}$)}
              \item{${B_{0}}$ is the magnetic flux density in region 0 (${\text{T}}$)}
              \end{symbDescription}
\\ \midrule
Notes & Combining Newton's second law with the Lorentz force law, (\hyperref[TM:lorentzForce]{TM:lorentzForce}), the equations of motion for each velocity component are obtained. The quantity q/m is the charge-to-mass ratio. Assumes \hyperref[singleParticle]{A:singleParticle}, \hyperref[prescribedFields]{A:prescribedFields}, \hyperref[lorentzOnly]{A:lorentzOnly}.
        
\\ \midrule
Source & --
         
\\ \midrule
RefBy & \hyperref[GD:dyn2D]{GD:dyn2D}
        
\\ \bottomrule
\end{tabular}
\end{minipage}

\subsubsection{General Definitions}
\label{Sec:GDs}
This section collects the laws and equations that will be used to build the instance models.

\medskip
\noindent
\begin{minipage}{\textwidth}
\begin{tabular}{>{\raggedright}p{0.13\textwidth}>{\raggedright\arraybackslash}p{0.82\textwidth}}
\toprule \textbf{Refname} & \textbf{GD:kin2D}
\phantomsection 
\label{GD:kin2D}
\\ \midrule
Label & Two-dimensional kinematics in Cartesian coordinates
        
\\ \midrule
Units & $\frac{\text{m}}{\text{s}}$
        
\\ \midrule
Equation & \begin{displaymath}
           \resizeExpression{\frac{\,dx}{\,dt}={v_{\text{x}}}\land{}\frac{\,dy}{\,dt}={v_{\text{y}}}}
           \end{displaymath}
\\ \midrule
Description & \begin{symbDescription}
              \item{$t$ is the time (${\text{s}}$)}
              \item{$x$ is the x-position of the particle (${\text{m}}$)}
              \item{${v_{\text{x}}}$ is the x-velocity of the particle ($\frac{\text{m}}{\text{s}}$)}
              \item{$y$ is the y-position of the particle (${\text{m}}$)}
              \item{${v_{\text{y}}}$ is the y-velocity of the particle ($\frac{\text{m}}{\text{s}}$)}
              \end{symbDescription}
\\ \midrule
Notes & The planar motion assumption (\hyperref[twoDMotion]{A:twoDMotion}) allows the particle position to be represented by $x$ and $y$ with velocity components ${v_{\text{x}}}$ and ${v_{\text{y}}}$ respectively. Velocity is defined in \hyperref[TM:velocity]{TM:velocity}..
        
\\ \midrule
Source & --
         
\\ \midrule
RefBy & \hyperref[IM:stateEvol]{IM:stateEvol}
        
\\ \bottomrule
\end{tabular}
\end{minipage}

\medskip
\noindent
\begin{minipage}{\textwidth}
\begin{tabular}{>{\raggedright}p{0.13\textwidth}>{\raggedright\arraybackslash}p{0.82\textwidth}}
\toprule \textbf{Refname} & \textbf{GD:vCrossB2D}
\phantomsection 
\label{GD:vCrossB2D}
\\ \midrule
Label & Cross product v x B for out-of-plane magnetic field
        
\\ \midrule
Units & $\frac{\text{N}}{\text{C}}$
        
\\ \midrule
Equation & \begin{displaymath}
           \resizeExpression{\symbf{v}\text{×}\symbf{B}=\begin{bmatrix}
                                                        {v_{\text{y}}}\,{B_{0}} & -{v_{\text{x}}}\,{B_{0}} & 0
                                                        \end{bmatrix}}
           \end{displaymath}
\\ \midrule
Description & \begin{symbDescription}
              \item{$\symbf{v}\text{×}\symbf{B}$ is the cross product of velocity and magnetic field (Unitless)}
              \item{${v_{\text{y}}}$ is the y-velocity of the particle ($\frac{\text{m}}{\text{s}}$)}
              \item{${B_{0}}$ is the magnetic flux density in region 0 (${\text{T}}$)}
              \item{${v_{\text{x}}}$ is the x-velocity of the particle ($\frac{\text{m}}{\text{s}}$)}
              \end{symbDescription}
\\ \midrule
Notes & Under planar motion (\hyperref[twoDMotion]{A:twoDMotion}) and an out-of-plane magnetic field (\hyperref[bPerpPlane]{A:bPerpPlane}), the velocity has the form v(t) = (vx(t), vy(t), 0) and the magnetic field has the form B = (0, 0, B). The cross product then reduces to the stated component form..
        
\\ \midrule
Source & --
         
\\ \midrule
RefBy & \hyperref[GD:dyn2D]{GD:dyn2D}
        
\\ \bottomrule
\end{tabular}
\end{minipage}

\medskip
\noindent
\begin{minipage}{\textwidth}
\begin{tabular}{>{\raggedright}p{0.13\textwidth}>{\raggedright\arraybackslash}p{0.82\textwidth}}
\toprule \textbf{Refname} & \textbf{GD:dyn2D}
\phantomsection 
\label{GD:dyn2D}
\\ \midrule
Label & Two-dimensional equations of motion under E and B
        
\\ \midrule
Units & $\frac{\text{m}}{\text{s}^{2}}$
        
\\ \midrule
Equation & \begin{displaymath}
           \resizeExpression{\frac{\,d{v_{\text{x}}}}{\,dt}=κ\,\left({E_{\text{x}0}}+{v_{\text{y}}}\,{B_{0}}\right)\land{}\frac{\,d{v_{\text{y}}}}{\,dt}=κ\,\left({E_{\text{y}0}}-{v_{\text{x}}}\,{B_{0}}\right)}
           \end{displaymath}
\\ \midrule
Description & \begin{symbDescription}
              \item{$t$ is the time (${\text{s}}$)}
              \item{${v_{\text{x}}}$ is the x-velocity of the particle ($\frac{\text{m}}{\text{s}}$)}
              \item{$κ$ is the charge-to-mass ratio ($\frac{\text{C}}{\text{kg}}$)}
              \item{${E_{\text{x}0}}$ is the x-electric field in region 0 ($\frac{\text{N}}{\text{C}}$)}
              \item{${v_{\text{y}}}$ is the y-velocity of the particle ($\frac{\text{m}}{\text{s}}$)}
              \item{${B_{0}}$ is the magnetic flux density in region 0 (${\text{T}}$)}
              \item{${E_{\text{y}0}}$ is the y-electric field in region 0 ($\frac{\text{N}}{\text{C}}$)}
              \end{symbDescription}
\\ \midrule
Notes & Assuming planar motion (\hyperref[twoDMotion]{A:twoDMotion}), prescribed fields (\hyperref[prescribedFields]{A:prescribedFields}), out-of-plane magnetic field (\hyperref[bPerpPlane]{A:bPerpPlane}), and axis-aligned electric field (\hyperref[eAxisAligned]{A:eAxisAligned}). Here $κ$ = q/m is the charge-to-mass ratio defined in \hyperref[DD:qOverm]{DD:qOverm}. The equations are obtained from \hyperref[TM:lorentzForce]{TM:lorentzForce} and \hyperref[TM:eqnMotion]{TM:eqnMotion} by applying \hyperref[GD:vCrossB2D]{GD:vCrossB2D}.
        
\\ \midrule
Source & --
         
\\ \midrule
RefBy & \hyperref[IM:stateEvol]{IM:stateEvol}
        
\\ \bottomrule
\end{tabular}
\end{minipage}

\subsubsection{Data Definitions}
\label{Sec:DDs}
This section collects and defines all the data needed to build the instance models.

\medskip
\noindent
\begin{minipage}{\textwidth}
\begin{tabular}{>{\raggedright}p{0.13\textwidth}>{\raggedright\arraybackslash}p{0.82\textwidth}}
\toprule \textbf{Refname} & \textbf{DD:qOverm}
\phantomsection 
\label{DD:qOverm}
\\ \midrule
Label & Charge-to-mass ratio
        
\\ \midrule
Symbol & $κ$
         
\\ \midrule
Units & $\frac{\text{C}}{\text{kg}}$
        
\\ \midrule
Equation & \begin{displaymath}
           \resizeExpression{κ=\frac{q}{m}}
           \end{displaymath}
\\ \midrule
Description & \begin{symbDescription}
              \item{$κ$ is the charge-to-mass ratio ($\frac{\text{C}}{\text{kg}}$)}
              \item{$q$ is the particle charge (${\text{C}}$)}
              \item{$m$ is the particle mass (${\text{kg}}$)}
              \end{symbDescription}
\\ \midrule
Notes & $q$ is the particle charge and $m$ is the particle mass. The ratio $κ$ is used to simplify the equations of motion.
        
\\ \midrule
Source & --
         
\\ \midrule
RefBy & \hyperref[IM:stateEvol]{IM:stateEvol} and \hyperref[GD:dyn2D]{GD:dyn2D}
        
\\ \bottomrule
\end{tabular}
\end{minipage}

\medskip
\noindent
\begin{minipage}{\textwidth}
\begin{tabular}{>{\raggedright}p{0.13\textwidth}>{\raggedright\arraybackslash}p{0.82\textwidth}}
\toprule \textbf{Refname} & \textbf{DD:initState}
\phantomsection 
\label{DD:initState}
\\ \midrule
Label & Initial state vector
        
\\ \midrule
Symbol & ${s_{0}}$
         
\\ \midrule
Units & Unitless
        
\\ \midrule
Equation & \begin{displaymath}
           \resizeExpression{{s_{0}}=\begin{bmatrix}
                                     {x_{0}} & {y_{0}} & {v_{\text{x}0}} & {v_{\text{y}0}}
                                     \end{bmatrix}}
           \end{displaymath}
\\ \midrule
Description & \begin{symbDescription}
              \item{${s_{0}}$ is the initial state vector (Unitless)}
              \item{${x_{0}}$ is the initial x-position (${\text{m}}$)}
              \item{${y_{0}}$ is the initial y-position (${\text{m}}$)}
              \item{${v_{\text{x}0}}$ is the initial x-velocity ($\frac{\text{m}}{\text{s}}$)}
              \item{${v_{\text{y}0}}$ is the initial y-velocity ($\frac{\text{m}}{\text{s}}$)}
              \end{symbDescription}
\\ \midrule
Notes & ${x_{0}}$, ${y_{0}}$ are the initial position components (m) and ${v_{\text{x}0}}$, ${v_{\text{y}0}}$ are the initial velocity components (m/s).
        
\\ \midrule
Source & --
         
\\ \midrule
RefBy & \hyperref[IM:stateEvol]{IM:stateEvol}
        
\\ \bottomrule
\end{tabular}
\end{minipage}

\medskip
\noindent
\begin{minipage}{\textwidth}
\begin{tabular}{>{\raggedright}p{0.13\textwidth}>{\raggedright\arraybackslash}p{0.82\textwidth}}
\toprule \textbf{Refname} & \textbf{DD:eField}
\phantomsection 
\label{DD:eField}
\\ \midrule
Label & Electric field vector
        
\\ \midrule
Symbol & $\symbf{E}$
         
\\ \midrule
Units & $\frac{\text{N}}{\text{C}}$
        
\\ \midrule
Equation & \begin{displaymath}
           \resizeExpression{\symbf{E}=\begin{bmatrix}
                                       {E_{\text{x}0}} & {E_{\text{y}0}} & 0
                                       \end{bmatrix}}
           \end{displaymath}
\\ \midrule
Description & \begin{symbDescription}
              \item{$\symbf{E}$ is the electric field vector ($\frac{\text{N}}{\text{C}}$)}
              \item{${E_{\text{x}0}}$ is the x-electric field in region 0 ($\frac{\text{N}}{\text{C}}$)}
              \item{${E_{\text{y}0}}$ is the y-electric field in region 0 ($\frac{\text{N}}{\text{C}}$)}
              \end{symbDescription}
\\ \midrule
Notes & $\symbf{E}$ is the electric field vector. ${E_{\text{x}0}}$ and ${E_{\text{y}0}}$ are the in-plane components (N/C). The out-of-plane component is zero by \hyperref[eAxisAligned]{A:eAxisAligned}, and the fields are user-specified per \hyperref[prescribedFields]{A:prescribedFields}.
        
\\ \midrule
Source & --
         
\\ \midrule
RefBy & 
\\ \bottomrule
\end{tabular}
\end{minipage}

\medskip
\noindent
\begin{minipage}{\textwidth}
\begin{tabular}{>{\raggedright}p{0.13\textwidth}>{\raggedright\arraybackslash}p{0.82\textwidth}}
\toprule \textbf{Refname} & \textbf{DD:bField}
\phantomsection 
\label{DD:bField}
\\ \midrule
Label & Magnetic flux density vector
        
\\ \midrule
Symbol & $\symbf{B}$
         
\\ \midrule
Units & ${\text{T}}$
        
\\ \midrule
Equation & \begin{displaymath}
           \resizeExpression{\symbf{B}=\begin{bmatrix}
                                       0 & 0 & {B_{0}}
                                       \end{bmatrix}}
           \end{displaymath}
\\ \midrule
Description & \begin{symbDescription}
              \item{$\symbf{B}$ is the magnetic flux density vector (${\text{T}}$)}
              \item{${B_{0}}$ is the magnetic flux density in region 0 (${\text{T}}$)}
              \end{symbDescription}
\\ \midrule
Notes & $\symbf{B}$ is the magnetic flux density vector. ${B_{0}}$ is the out-of-plane (z) component (T). The in-plane components are zero by \hyperref[bPerpPlane]{A:bPerpPlane}; fields are fixed by \hyperref[prescribedFields]{A:prescribedFields}.
        
\\ \midrule
Source & --
         
\\ \midrule
RefBy & 
\\ \bottomrule
\end{tabular}
\end{minipage}

\medskip
\noindent
\begin{minipage}{\textwidth}
\begin{tabular}{>{\raggedright}p{0.13\textwidth}>{\raggedright\arraybackslash}p{0.82\textwidth}}
\toprule \textbf{Refname} & \textbf{DD:regionRect}
\phantomsection 
\label{DD:regionRect}
\\ \midrule
Label & Rectangular field region
        
\\ \midrule
Symbol & ${R_{i}}$
         
\\ \midrule
Units & Unitless
        
\\ \midrule
Equation & \begin{displaymath}
           \resizeExpression{{R_{i}}=\begin{bmatrix}
                                     {x_{\text{grid}}} & {y_{\text{grid}}} & w & h
                                     \end{bmatrix}}
           \end{displaymath}
\\ \midrule
Description & \begin{symbDescription}
              \item{${R_{i}}$ is the rectangular field region (Unitless)}
              \item{${x_{\text{grid}}}$ is the grid origin x-coordinate (${\text{m}}$)}
              \item{${y_{\text{grid}}}$ is the grid origin y-coordinate (${\text{m}}$)}
              \item{$w$ is the region width (${\text{m}}$)}
              \item{$h$ is the region height (${\text{m}}$)}
              \end{symbDescription}
\\ \midrule
Notes & Each field region R\_i (for i = 1 .. $N$ ) is an axis-aligned rectangle of width $w$ and height $h$. The equation above lists the grid parameters that define each region. For a row layout (regions tiled left-to-right): R\_i = [x\_grid + (i-1)*w, x\_grid + i*w] x [y\_grid, y\_grid + h]. For a column layout (regions tiled bottom-to-top): R\_i = [x\_grid, x\_grid + w] x [y\_grid + (i-1)*h, y\_grid + i*h]. The grid origin (${x_{\text{grid}}}$, ${y_{\text{grid}}}$) anchors region 1, per \hyperref[rectRegions]{A:rectRegions} and \hyperref[gridLayout]{A:gridLayout}.
        
\\ \midrule
Source & --
         
\\ \midrule
RefBy & \hyperref[buildRegionGrid]{FR:Build-Region-Grid} and \hyperref[DD:fieldsByRegion]{DD:fieldsByRegion}
        
\\ \bottomrule
\end{tabular}
\end{minipage}

\medskip
\noindent
\begin{minipage}{\textwidth}
\begin{tabular}{>{\raggedright}p{0.13\textwidth}>{\raggedright\arraybackslash}p{0.82\textwidth}}
\toprule \textbf{Refname} & \textbf{DD:fieldsByRegion}
\phantomsection 
\label{DD:fieldsByRegion}
\\ \midrule
Label & Electric field in region i
        
\\ \midrule
Symbol & ${\symbf{E}_{i}}$
         
\\ \midrule
Units & $\frac{\text{N}}{\text{C}}$
        
\\ \midrule
Equation & \begin{displaymath}
           \resizeExpression{{\symbf{E}_{i}}=\begin{bmatrix}
                                             {E_{\text{x}0}} & {E_{\text{y}0}} & {B_{0}}
                                             \end{bmatrix}}
           \end{displaymath}
\\ \midrule
Description & \begin{symbDescription}
              \item{${\symbf{E}_{i}}$ is the electric field in region i ($\frac{\text{N}}{\text{C}}$)}
              \item{${E_{\text{x}0}}$ is the x-electric field in region 0 ($\frac{\text{N}}{\text{C}}$)}
              \item{${E_{\text{y}0}}$ is the y-electric field in region 0 ($\frac{\text{N}}{\text{C}}$)}
              \item{${B_{0}}$ is the magnetic flux density in region 0 (${\text{T}}$)}
              \end{symbDescription}
\\ \midrule
Notes & Each region R\_i has its own independently specified electric field ${\symbf{E}_{i}}$ = (Ex\_i, Ey\_i, 0) and magnetic flux density ${\symbf{B}_{i}}$ = (0, 0, B\_i). The fields are constant within each region per \hyperref[piecewiseUniform]{A:piecewiseUniform}. The active fields at time t are determined by which region R\_i (\hyperref[DD:regionRect]{DD:regionRect}) contains the particle position (x(t), y(t)). If the particle is outside all defined regions, the fields default to zero and the particle undergoes free (force-free) motion. The per-region values (Ex\_i, Ey\_i, B\_i) are user-specified inputs per \hyperref[prescribedFields]{A:prescribedFields}.
        
\\ \midrule
Source & --
         
\\ \midrule
RefBy & \hyperref[IM:stateEvol]{IM:stateEvol} and \hyperref[lookupActiveRegion]{FR:Lookup-Active-Region}
        
\\ \bottomrule
\end{tabular}
\end{minipage}

\medskip
\noindent
\begin{minipage}{\textwidth}
\begin{tabular}{>{\raggedright}p{0.13\textwidth}>{\raggedright\arraybackslash}p{0.82\textwidth}}
\toprule \textbf{Refname} & \textbf{DD:detectorLine}
\phantomsection 
\label{DD:detectorLine}
\\ \midrule
Label & Detector line
        
\\ \midrule
Symbol & ${L_{\text{det}}}$
         
\\ \midrule
Units & Unitless
        
\\ \midrule
Equation & \begin{displaymath}
           \resizeExpression{{L_{\text{det}}}=\begin{bmatrix}
                                              {d_{\text{orient}}} & {x_{\text{det}}} & {y_{\text{det}}} & {y_{\text{detMin}}} & {y_{\text{detMax}}} & {x_{\text{detMin}}} & {x_{\text{detMax}}}
                                              \end{bmatrix}}
           \end{displaymath}
\\ \midrule
Description & \begin{symbDescription}
              \item{${L_{\text{det}}}$ is the detector line (Unitless)}
              \item{${d_{\text{orient}}}$ is the detector orientation (Unitless)}
              \item{${x_{\text{det}}}$ is the detector line x-position (${\text{m}}$)}
              \item{${y_{\text{det}}}$ is the detector line y-position (${\text{m}}$)}
              \item{${y_{\text{detMin}}}$ is the minimum y-coordinate of detector (${\text{m}}$)}
              \item{${y_{\text{detMax}}}$ is the maximum y-coordinate of detector (${\text{m}}$)}
              \item{${x_{\text{detMin}}}$ is the minimum x-coordinate of detector (${\text{m}}$)}
              \item{${x_{\text{detMax}}}$ is the maximum x-coordinate of detector (${\text{m}}$)}
              \end{symbDescription}
\\ \midrule
Notes & ${L_{\text{det}}}$ is the detector line segment. When ${d_{\text{orient}}}$ = 0 (vertical), the detector is L\_det is the set of (x,y) such that x = x\_det and y in [y\_min(det), y\_max(det)]. When ${d_{\text{orient}}}$ = 1 (horizontal), the detector is L\_det is the set of (x,y) such that y = y\_det and x in [x\_min(det), x\_max(det)]. The detector orientation and bounds are user-specified, per \hyperref[lineDetector]{A:lineDetector}.
        
\\ \midrule
Source & --
         
\\ \midrule
RefBy & \hyperref[IM:detHit]{IM:detHit}
        
\\ \bottomrule
\end{tabular}
\end{minipage}

\subsubsection{Instance Models}
\label{Sec:IMs}
This section transforms the problem defined in the \hyperref[Sec:ProbDesc]{problem description} into one which is expressed in mathematical terms. It uses concrete symbols defined in the \hyperref[Sec:DDs]{data definitions} to replace the abstract symbols in the models identified in \hyperref[Sec:TMs]{theoretical models} and \hyperref[Sec:GDs]{general definitions}.

\medskip
\noindent
\begin{minipage}{\textwidth}
\begin{tabular}{>{\raggedright}p{0.13\textwidth}>{\raggedright\arraybackslash}p{0.82\textwidth}}
\toprule \textbf{Refname} & \textbf{IM:stateEvol}
\phantomsection 
\label{IM:stateEvol}
\\ \midrule
Label & Particle state evolution in the x-y plane
        
\\ \midrule
Input & $m$, $q$, ${x_{0}}$, ${y_{0}}$, ${v_{\text{x}0}}$, ${v_{\text{y}0}}$, ${E_{\text{x}0}}$, ${E_{\text{y}0}}$, ${B_{0}}$, ${t_{\text{final}}}$
        
\\ \midrule
Output & $x$
         
\\ \midrule
Input Constraints & 
\\ \midrule
Output Constraints & 
\\ \midrule
Equation & \begin{displaymath}
           \resizeExpression{\frac{\,dx}{\,dt}={v_{\text{x}}}\land{}\frac{\,dy}{\,dt}={v_{\text{y}}}\land{}\frac{\,d{v_{\text{x}}}}{\,dt}=κ\,\left({E_{\text{x}0}}+{v_{\text{y}}}\,{B_{0}}\right)\land{}\frac{\,d{v_{\text{y}}}}{\,dt}=κ\,\left({E_{\text{y}0}}-{v_{\text{x}}}\,{B_{0}}\right)}
           \end{displaymath}
\\ \midrule
Description & \begin{symbDescription}
              \item{$t$ is the time (${\text{s}}$)}
              \item{$x$ is the x-position of the particle (${\text{m}}$)}
              \item{${v_{\text{x}}}$ is the x-velocity of the particle ($\frac{\text{m}}{\text{s}}$)}
              \item{$y$ is the y-position of the particle (${\text{m}}$)}
              \item{${v_{\text{y}}}$ is the y-velocity of the particle ($\frac{\text{m}}{\text{s}}$)}
              \item{$κ$ is the charge-to-mass ratio ($\frac{\text{C}}{\text{kg}}$)}
              \item{${E_{\text{x}0}}$ is the x-electric field in region 0 ($\frac{\text{N}}{\text{C}}$)}
              \item{${B_{0}}$ is the magnetic flux density in region 0 (${\text{T}}$)}
              \item{${E_{\text{y}0}}$ is the y-electric field in region 0 ($\frac{\text{N}}{\text{C}}$)}
              \end{symbDescription}
\\ \midrule
Notes & Bold symbols denote vectors; the state vector s(t) = (x(t), y(t), vx(t), vy(t)) collects position and velocity components. $κ$ = q/m is the charge-to-mass ratio (\hyperref[DD:qOverm]{DD:qOverm}). The initial condition is s(0) = (x0, y0, v0x, v0y) (\hyperref[DD:initState]{DD:initState}). The fields (Ex(t), Ey(t)) and B(t) are determined by which region R\_i contains (x(t), y(t)) at time t (\hyperref[DD:fieldsByRegion]{DD:fieldsByRegion}). If (x(t), y(t)) lies outside all defined regions, the fields default to zero and the particle undergoes free (force-free) motion. The model is derived by combining \hyperref[GD:kin2D]{GD:kin2D} and \hyperref[GD:dyn2D]{GD:dyn2D}. Applicable assumptions: \hyperref[singleParticle]{A:singleParticle}, \hyperref[noInteractions]{A:noInteractions}, \hyperref[twoDMotion]{A:twoDMotion} and \hyperref[lorentzOnly]{A:lorentzOnly}.
        
\\ \midrule
Source & --
         
\\ \midrule
RefBy & \hyperref[IM:detHit]{IM:detHit}, \hyperref[reportOutputs]{FR:Report-Outputs}, \hyperref[predictTrajectory]{GS:predictTrajectory}, and \hyperref[computeTraj]{FR:Compute-Trajectory}
        
\\ \bottomrule
\end{tabular}
\end{minipage}

\medskip
\noindent
\begin{minipage}{\textwidth}
\begin{tabular}{>{\raggedright}p{0.13\textwidth}>{\raggedright\arraybackslash}p{0.82\textwidth}}
\toprule \textbf{Refname} & \textbf{IM:detHit}
\phantomsection 
\label{IM:detHit}
\\ \midrule
Label & First detector intersection (hit time and hit location)
        
\\ \midrule
Input & ${d_{\text{orient}}}$, ${x_{\text{det}}}$, ${y_{\text{det}}}$, ${y_{\text{detMin}}}$, ${y_{\text{detMax}}}$, ${x_{\text{detMin}}}$, ${x_{\text{detMax}}}$, ${t_{\text{final}}}$, $x$, $y$
        
\\ \midrule
Output & ${t_{\text{hit}}}$
         
\\ \midrule
Input Constraints & 
\\ \midrule
Output Constraints & 
\\ \midrule
Equation & \begin{displaymath}
           \resizeExpression{x={x_{\text{det}}}\land{}{y_{\text{hit}}}\geq{}{y_{\text{detMin}}}\land{}{y_{\text{hit}}}\leq{}{y_{\text{detMax}}}\lor{}y={y_{\text{det}}}\land{}{x_{\text{hit}}}\geq{}{x_{\text{detMin}}}\land{}{x_{\text{hit}}}\leq{}{x_{\text{detMax}}}}
           \end{displaymath}
\\ \midrule
Description & \begin{symbDescription}
              \item{$x$ is the x-position of the particle (${\text{m}}$)}
              \item{${x_{\text{det}}}$ is the detector line x-position (${\text{m}}$)}
              \item{${y_{\text{hit}}}$ is the y-coordinate of impact point (${\text{m}}$)}
              \item{${y_{\text{detMin}}}$ is the minimum y-coordinate of detector (${\text{m}}$)}
              \item{${y_{\text{detMax}}}$ is the maximum y-coordinate of detector (${\text{m}}$)}
              \item{$y$ is the y-position of the particle (${\text{m}}$)}
              \item{${y_{\text{det}}}$ is the detector line y-position (${\text{m}}$)}
              \item{${x_{\text{hit}}}$ is the x-coordinate of impact point (${\text{m}}$)}
              \item{${x_{\text{detMin}}}$ is the minimum x-coordinate of detector (${\text{m}}$)}
              \item{${x_{\text{detMax}}}$ is the maximum x-coordinate of detector (${\text{m}}$)}
              \end{symbDescription}
\\ \midrule
Notes & The detector is modelled as a line segment that may be vertical or horizontal, selected by the orientation flag ${d_{\text{orient}}}$ (\hyperref[DD:detectorLine]{DD:detectorLine}). Vertical case (D = 0): the detector is at x = ${x_{\text{det}}}$ with y-bounds [${y_{\text{detMin}}}$, ${y_{\text{detMax}}}$]. The hit condition is x(t) = x\_det and y(t) in [y\_min, y\_max]. Horizontal case (D = 1): the detector is at y = ${y_{\text{det}}}$ with x-bounds [${x_{\text{detMin}}}$, ${x_{\text{detMax}}}$]. The hit condition is y(t) = y\_det and x(t) in [x\_min, x\_max]. In both cases, ${t_{\text{hit}}}$ is the minimum t in [0, ${t_{\text{final}}}$] satisfying the condition. If no such t exists, ${t_{\text{hit}}}$ is set to -1 as a sentinel value indicating no hit, and the hit coordinates are undefined. The trajectory (x(t), y(t)) is provided by \hyperref[IM:stateEvol]{IM:stateEvol}. Applicable assumptions: \hyperref[fullDetection]{A:fullDetection} and \hyperref[lineDetector]{A:lineDetector}.
        
\\ \midrule
Source & --
         
\\ \midrule
RefBy & \hyperref[reportOutputs]{FR:Report-Outputs} and \hyperref[determineDetectorOutcome]{GS:determineDetectorOutcome}
        
\\ \bottomrule
\end{tabular}
\end{minipage}

\subsubsection{Data Constraints}
\label{Sec:DataConstraints}
The \hyperref[Table:InDataConstraints]{Input Data Constraints Table} shows the data constraints on the input variables. The column for physical constraints gives the physical limitations on the range of values that can be taken by the variable. The uncertainty column provides an estimate of the confidence with which the physical quantities can be measured. This information would be part of the input if one were performing an uncertainty quantification exercise. The constraints are conservative to give the user of the model the flexibility to experiment with unusual situations. The column of typical values is intended to provide a feel for a common scenario.

\begin{longtblr}
[caption={Input Data Constraints}]
{colspec={l l l l l}, rowhead=1, hline{1,Z}=\heavyrulewidth, hline{2}=\lightrulewidth}
\textbf{Var} & \textbf{Physical Constraints} & \textbf{Software Constraints} & \textbf{Typical Value} & \textbf{Uncert.}
\\
${d_{\text{orient}}}$ & -- & $0\leq{}{d_{\text{orient}}}\leq{}1$ & $0$ & 0.0$\%$
\\
$h$ & $h\gt{}0$ & -- & $0.1$ ${\text{m}}$ & 0.0$\%$
\\
$m$ & $m\gt{}0$ & ${m_{\text{min}}}\leq{}m\leq{}{m_{\text{max}}}$ & $911.0\cdot{}10^{-33}$ ${\text{kg}}$ & 10.0$\%$
\\
$N$ & $N\gt{}0$ & -- & $1$ & 0.0$\%$
\\
${N_{\text{col}}}$ & ${N_{\text{col}}}\gt{}0$ & -- & $1$ & 0.0$\%$
\\
$q$ & -- & $-{q_{\text{max}}}\leq{}q\leq{}{q_{\text{max}}}$ & $160.0\cdot{}10^{-21}$ ${\text{C}}$ & 10.0$\%$
\\
${t_{\text{final}}}$ & ${t_{\text{final}}}\gt{}0$ & ${t_{\text{final}}}\leq{}{t_{\text{max}}}$ & $1.0\cdot{}10^{-6}$ ${\text{s}}$ & 0.0$\%$
\\
${v_{\text{x}0}}$ & -- & $-{v_{\text{max}}}\leq{}{v_{\text{x}0}}\leq{}{v_{\text{max}}}$ & $1.0\cdot{}10^{6}$ $\frac{\text{m}}{\text{s}}$ & 10.0$\%$
\\
${v_{\text{y}0}}$ & -- & $-{v_{\text{max}}}\leq{}{v_{\text{y}0}}\leq{}{v_{\text{max}}}$ & $0$ $\frac{\text{m}}{\text{s}}$ & 10.0$\%$
\\
$w$ & $w\gt{}0$ & -- & $0.1$ ${\text{m}}$ & 0.0$\%$
\\
${x_{0}}$ & -- & -- & $0$ ${\text{m}}$ & 0.0$\%$
\\
${x_{\text{det}}}$ & -- & -- & $0.1$ ${\text{m}}$ & 0.0$\%$
\\
${x_{\text{detMax}}}$ & -- & -- & $0.1$ ${\text{m}}$ & 0.0$\%$
\\
${x_{\text{detMin}}}$ & -- & -- & $0$ ${\text{m}}$ & 0.0$\%$
\\
${x_{\text{grid}}}$ & -- & -- & $0$ ${\text{m}}$ & 0.0$\%$
\\
${y_{0}}$ & -- & -- & $0$ ${\text{m}}$ & 0.0$\%$
\\
${y_{\text{det}}}$ & -- & -- & $0$ ${\text{m}}$ & 0.0$\%$
\\
${y_{\text{detMax}}}$ & -- & -- & $0.05$ ${\text{m}}$ & 0.0$\%$
\\
${y_{\text{detMin}}}$ & -- & -- & $0.0-5$ ${\text{m}}$ & 0.0$\%$
\\
${y_{\text{grid}}}$ & -- & -- & $0$ ${\text{m}}$ & 0.0$\%$
\label{Table:InDataConstraints}
\end{longtblr}
\subsubsection{Properties of a Correct Solution}
\label{Sec:CorSolProps}
The \hyperref[Table:OutDataConstraints]{Output Data Constraints Table} shows the data constraints on the output variables. The column for physical constraints gives the physical limitations on the range of values that can be taken by the variable.

\begin{longtblr}
[caption={Output Data Constraints}]
{colspec={l}, rowhead=1, hline{1,Z}=\heavyrulewidth, hline{2}=\lightrulewidth}
\textbf{Var}
\\
$s$
\\
${t_{\text{hit}}}$
\\
${y_{\text{hit}}}$
\\
${x_{\text{hit}}}$
\label{Table:OutDataConstraints}
\end{longtblr}
For a pure magnetic field (no electric field, i.e., Ex = Ey = 0), the speed of the particle must remain constant over time, since the magnetic force is always perpendicular to the velocity and therefore does no work on the particle. A correct solution must satisfy |v(t)| = |v(0)| for all times t in the simulation.

When the electric field is non-zero, the work-energy theorem requires that the change in kinetic energy equals the work done by the electric force: delta\_KE = q * Ex * delta\_x + q * Ey * delta\_y. A correct solution must respect this energy balance over each time step.

In the absence of all fields (Ex = Ey = 0 and B = 0), Newton's first law requires that the particle travel in a straight line at constant velocity. A correct solution must show no change in either velocity component and linear growth in both position components.

When the particle crosses from region R\_i to region R\_j, the position and velocity must be continuous at the boundary. Only the acceleration changes (due to the different fields in the new region). A correct solution must preserve continuity of the state vector at every region boundary crossing..

\section{Requirements}
\label{Sec:Requirements}
This section provides the functional requirements, the tasks and behaviours that the software is expected to complete, and the non-functional requirements, the qualities that the software is expected to exhibit.

\subsection{Functional Requirements}
\label{Sec:FRs}
This section provides the functional requirements, the tasks and behaviours that the software is expected to complete.

\begin{description}[font=\normalfont]
\item[Input-Values:\phantomsection\label{inputValues}]{Input the values from \hyperref[Table:ReqInputs]{Tab:ReqInputs}.}
\item[Echo-Inputs:\phantomsection\label{echoInputs}]{Echo the given inputs.}
\item[Build-Region-Grid:\phantomsection\label{buildRegionGrid}]{Construct the N field regions from the grid origin, region dimensions (w, h), and region count N, and assign each region its specified (Ex\_i, Ey\_i, B\_i) (\hyperref[DD:regionRect]{DD:regionRect}).}
\item[Lookup-Active-Region:\phantomsection\label{lookupActiveRegion}]{At each simulation time-step, determine which region R\_i contains the particle position (x(t), y(t)) and apply the corresponding fields. If the particle is outside all regions, apply zero fields (\hyperref[DD:fieldsByRegion]{DD:fieldsByRegion}).}
\item[Compute-Trajectory:\phantomsection\label{computeTraj}]{Compute the trajectory of the charged particle using the equations of motion derived from the Lorentz force (\hyperref[IM:stateEvol]{IM:stateEvol}).}
\item[Report-Outputs:\phantomsection\label{reportOutputs}]{Report the final position, velocity, and whether the particle reaches the detector line, based on \hyperref[IM:stateEvol]{IM:stateEvol} and \hyperref[IM:detHit]{IM:detHit}.}
\end{description}
\begin{longtblr}
[caption={Required Inputs}]
{colspec={l l l}, rowhead=1, hline{1,Z}=\heavyrulewidth, hline{2}=\lightrulewidth}
\textbf{Symbol} & \textbf{Description} & \textbf{Units}
\\
${B_{0}}$ & Magnetic flux density in region 0 & ${\text{T}}$
\\
${B_{1}}$ & Magnetic flux density in region 1 & ${\text{T}}$
\\
${B_{2}}$ & Magnetic flux density in region 2 & ${\text{T}}$
\\
${B_{3}}$ & Magnetic flux density in region 3 & ${\text{T}}$
\\
${B_{4}}$ & Magnetic flux density in region 4 & ${\text{T}}$
\\
${B_{5}}$ & Magnetic flux density in region 5 & ${\text{T}}$
\\
${d_{\text{orient}}}$ & Detector orientation & --
\\
${E_{\text{x}0}}$ & X-electric field in region 0 & $\frac{\text{N}}{\text{C}}$
\\
${E_{\text{x}1}}$ & X-electric field in region 1 & $\frac{\text{N}}{\text{C}}$
\\
${E_{\text{x}2}}$ & X-electric field in region 2 & $\frac{\text{N}}{\text{C}}$
\\
${E_{\text{x}3}}$ & X-electric field in region 3 & $\frac{\text{N}}{\text{C}}$
\\
${E_{\text{x}4}}$ & X-electric field in region 4 & $\frac{\text{N}}{\text{C}}$
\\
${E_{\text{x}5}}$ & X-electric field in region 5 & $\frac{\text{N}}{\text{C}}$
\\
${E_{\text{y}0}}$ & Y-electric field in region 0 & $\frac{\text{N}}{\text{C}}$
\\
${E_{\text{y}1}}$ & Y-electric field in region 1 & $\frac{\text{N}}{\text{C}}$
\\
${E_{\text{y}2}}$ & Y-electric field in region 2 & $\frac{\text{N}}{\text{C}}$
\\
${E_{\text{y}3}}$ & Y-electric field in region 3 & $\frac{\text{N}}{\text{C}}$
\\
${E_{\text{y}4}}$ & Y-electric field in region 4 & $\frac{\text{N}}{\text{C}}$
\\
${E_{\text{y}5}}$ & Y-electric field in region 5 & $\frac{\text{N}}{\text{C}}$
\\
$h$ & Region height & ${\text{m}}$
\\
$m$ & Particle mass & ${\text{kg}}$
\\
$N$ & Number of field regions & --
\\
${N_{\text{col}}}$ & Number of grid columns & --
\\
$q$ & Particle charge & ${\text{C}}$
\\
${t_{\text{final}}}$ & Final simulation time & ${\text{s}}$
\\
${v_{\text{x}0}}$ & Initial x-velocity & $\frac{\text{m}}{\text{s}}$
\\
${v_{\text{y}0}}$ & Initial y-velocity & $\frac{\text{m}}{\text{s}}$
\\
$w$ & Region width & ${\text{m}}$
\\
${x_{0}}$ & Initial x-position & ${\text{m}}$
\\
${x_{\text{det}}}$ & Detector line x-position & ${\text{m}}$
\\
${x_{\text{detMax}}}$ & Maximum x-coordinate of detector & ${\text{m}}$
\\
${x_{\text{detMin}}}$ & Minimum x-coordinate of detector & ${\text{m}}$
\\
${x_{\text{grid}}}$ & Grid origin x-coordinate & ${\text{m}}$
\\
${y_{0}}$ & Initial y-position & ${\text{m}}$
\\
${y_{\text{det}}}$ & Detector line y-position & ${\text{m}}$
\\
${y_{\text{detMax}}}$ & Maximum y-coordinate of detector & ${\text{m}}$
\\
${y_{\text{detMin}}}$ & Minimum y-coordinate of detector & ${\text{m}}$
\\
${y_{\text{grid}}}$ & Grid origin y-coordinate & ${\text{m}}$
\label{Table:ReqInputs}
\end{longtblr}
\subsection{Non-Functional Requirements}
\label{Sec:NFRs}
This section provides the non-functional requirements, the qualities that the software is expected to exhibit.

\begin{description}[font=\normalfont]
\item[Correctness:\phantomsection\label{correct}]{The outputs of the code have the \hyperref[Sec:CorSolProps]{properties of a correct solution}.}
\item[Maintainability:\phantomsection\label{maintainable}]{If a likely change is made to the finished software, it will take at most 10$\%$ of the original development time, assuming the same development resources are available.}
\item[Portability:\phantomsection\label{portable}]{The code shall be portable to multiple environments, particularly Windows, Mac OSX, and Linux.}
\end{description}
\section{Likely Changes}
\label{Sec:LCs}
This section lists the likely changes (LC) to be made to the software.

\begin{description}[font=\normalfont]
\item[lcExtendTo3D:\phantomsection\label{lcExtendTo3D}]{\hyperref[twoDMotion]{A:twoDMotion} - extend Trajecto to support three-dimensional (3D) motion, including a z-component of position, velocity, acceleration, electric field, and magnetic field.}
\end{description}
\section{Unlikely Changes}
\label{Sec:UCs}
This section lists the unlikely changes (UC) to be made to the software.

\begin{description}[font=\normalfont]
\item[ucLorentzForce:\phantomsection\label{ucLorentzForce}]{\hyperref[lorentzOnly]{A:lorentzOnly} - the Lorentz force law is a well-established physical principle and is unlikely to change.}
\end{description}
\section{Traceability Matrices and Graphs}
\label{Sec:TraceMatrices}
The purpose of the traceability matrices is to provide easy references on what has to be additionally modified if a certain component is changed. Every time a component is changed, the items in the column of that component that are marked with an ``X'' should be modified as well. \hyperref[Table:TraceMatAvsA]{Tab:TraceMatAvsA} shows the dependencies of the assumptions on each other. \hyperref[Table:TraceMatAvsAll]{Tab:TraceMatAvsAll} shows the dependencies of the data definitions, theoretical models, general definitions, instance models, requirements, likely changes, and unlikely changes on the assumptions. \hyperref[Table:TraceMatRefvsRef]{Tab:TraceMatRefvsRef} shows the dependencies of the data definitions, theoretical models, general definitions, and instance models on each other. \hyperref[Table:TraceMatAllvsR]{Tab:TraceMatAllvsR} shows the dependencies of the requirements and goal statements on the data definitions, theoretical models, general definitions, and instance models.

\begin{longtblr}
[caption={Traceability Matrix Showing the Connections Between Assumptions and Other Assumptions}]
{colspec={l l l l l l l l l l l l l}, rowhead=1, hline{1,Z}=\heavyrulewidth, hline{2}=\lightrulewidth}
\textbf{} & \textbf{\hyperref[piecewiseUniform]{A:piecewiseUniform}} & \textbf{\hyperref[singleParticle]{A:singleParticle}} & \textbf{\hyperref[noInteractions]{A:noInteractions}} & \textbf{\hyperref[prescribedFields]{A:prescribedFields}} & \textbf{\hyperref[twoDMotion]{A:twoDMotion}} & \textbf{\hyperref[bPerpPlane]{A:bPerpPlane}} & \textbf{\hyperref[eAxisAligned]{A:eAxisAligned}} & \textbf{\hyperref[rectRegions]{A:rectRegions}} & \textbf{\hyperref[gridLayout]{A:gridLayout}} & \textbf{\hyperref[lineDetector]{A:lineDetector}} & \textbf{\hyperref[fullDetection]{A:fullDetection}} & \textbf{\hyperref[lorentzOnly]{A:lorentzOnly}}
\\
\hyperref[piecewiseUniform]{A:piecewiseUniform} &  &  &  &  &  &  &  &  &  &  &  & 
\\
\hyperref[singleParticle]{A:singleParticle} &  &  &  &  &  &  &  &  &  &  &  & 
\\
\hyperref[noInteractions]{A:noInteractions} &  & X &  &  &  &  &  &  &  &  &  & 
\\
\hyperref[prescribedFields]{A:prescribedFields} &  &  &  &  &  &  &  &  &  &  &  & 
\\
\hyperref[twoDMotion]{A:twoDMotion} &  &  &  &  &  &  &  &  &  &  &  & 
\\
\hyperref[bPerpPlane]{A:bPerpPlane} &  &  &  &  & X &  &  &  &  &  &  & 
\\
\hyperref[eAxisAligned]{A:eAxisAligned} &  &  &  &  & X &  &  &  &  &  &  & 
\\
\hyperref[rectRegions]{A:rectRegions} &  &  &  &  &  &  &  &  &  &  &  & 
\\
\hyperref[gridLayout]{A:gridLayout} &  &  &  &  &  &  &  & X &  &  &  & 
\\
\hyperref[lineDetector]{A:lineDetector} &  &  &  &  &  &  &  &  &  &  &  & 
\\
\hyperref[fullDetection]{A:fullDetection} &  &  &  &  &  &  &  &  &  &  &  & 
\\
\hyperref[lorentzOnly]{A:lorentzOnly} &  &  &  &  &  &  &  &  &  &  &  & 
\label{Table:TraceMatAvsA}
\end{longtblr}
\begin{longtblr}
[caption={Traceability Matrix Showing the Connections Between Assumptions and Other Items}]
{colspec={l l l l l l l l l l l l l}, rowhead=1, hline{1,Z}=\heavyrulewidth, hline{2}=\lightrulewidth}
\textbf{} & \textbf{\hyperref[piecewiseUniform]{A:piecewiseUniform}} & \textbf{\hyperref[singleParticle]{A:singleParticle}} & \textbf{\hyperref[noInteractions]{A:noInteractions}} & \textbf{\hyperref[prescribedFields]{A:prescribedFields}} & \textbf{\hyperref[twoDMotion]{A:twoDMotion}} & \textbf{\hyperref[bPerpPlane]{A:bPerpPlane}} & \textbf{\hyperref[eAxisAligned]{A:eAxisAligned}} & \textbf{\hyperref[rectRegions]{A:rectRegions}} & \textbf{\hyperref[gridLayout]{A:gridLayout}} & \textbf{\hyperref[lineDetector]{A:lineDetector}} & \textbf{\hyperref[fullDetection]{A:fullDetection}} & \textbf{\hyperref[lorentzOnly]{A:lorentzOnly}}
\\
\hyperref[DD:qOverm]{DD:qOverm} &  &  &  &  &  &  &  &  &  &  &  & 
\\
\hyperref[DD:initState]{DD:initState} &  &  &  &  &  &  &  &  &  &  &  & 
\\
\hyperref[DD:eField]{DD:eField} &  &  &  & X &  &  & X &  &  &  &  & 
\\
\hyperref[DD:bField]{DD:bField} &  &  &  & X &  & X &  &  &  &  &  & 
\\
\hyperref[DD:regionRect]{DD:regionRect} &  &  &  &  &  &  &  & X & X &  &  & 
\\
\hyperref[DD:fieldsByRegion]{DD:fieldsByRegion} & X &  &  & X &  &  &  &  &  &  &  & 
\\
\hyperref[DD:detectorLine]{DD:detectorLine} &  &  &  &  &  &  &  &  &  & X &  & 
\\
\hyperref[TM:velocity]{TM:velocity} &  &  &  &  &  &  &  &  &  &  &  & 
\\
\hyperref[TM:acceleration]{TM:acceleration} &  &  &  &  &  &  &  &  &  &  &  & 
\\
\hyperref[TM:lorentzForce]{TM:lorentzForce} &  & X &  & X &  &  &  &  &  &  &  & 
\\
\hyperref[TM:eqnMotion]{TM:eqnMotion} &  & X &  & X &  &  &  &  &  &  &  & X
\\
\hyperref[GD:kin2D]{GD:kin2D} &  &  &  &  & X &  &  &  &  &  &  & 
\\
\hyperref[GD:vCrossB2D]{GD:vCrossB2D} &  &  &  &  & X & X &  &  &  &  &  & 
\\
\hyperref[GD:dyn2D]{GD:dyn2D} &  &  &  & X & X & X & X &  &  &  &  & 
\\
\hyperref[IM:stateEvol]{IM:stateEvol} &  & X & X &  & X &  &  &  &  &  &  & X
\\
\hyperref[IM:detHit]{IM:detHit} &  &  &  &  &  &  &  &  &  & X & X & 
\\
\hyperref[inputValues]{FR:Input-Values} &  &  &  &  &  &  &  &  &  &  &  & 
\\
\hyperref[echoInputs]{FR:Echo-Inputs} &  &  &  &  &  &  &  &  &  &  &  & 
\\
\hyperref[buildRegionGrid]{FR:Build-Region-Grid} &  &  &  &  &  &  &  &  &  &  &  & 
\\
\hyperref[lookupActiveRegion]{FR:Lookup-Active-Region} &  &  &  &  &  &  &  &  &  &  &  & 
\\
\hyperref[computeTraj]{FR:Compute-Trajectory} &  &  &  &  &  &  &  &  &  &  &  & 
\\
\hyperref[reportOutputs]{FR:Report-Outputs} &  &  &  &  &  &  &  &  &  &  &  & 
\\
\hyperref[correct]{NFR:Correctness} &  &  &  &  &  &  &  &  &  &  &  & 
\\
\hyperref[maintainable]{NFR:Maintainability} &  &  &  &  &  &  &  &  &  &  &  & 
\\
\hyperref[portable]{NFR:Portability} &  &  &  &  &  &  &  &  &  &  &  & 
\\
\hyperref[lcExtendTo3D]{LC:lcExtendTo3D} &  &  &  &  & X &  &  &  &  &  &  & 
\\
\hyperref[ucLorentzForce]{UC:ucLorentzForce} &  &  &  &  &  &  &  &  &  &  &  & X
\label{Table:TraceMatAvsAll}
\end{longtblr}
\begin{longtblr}
[caption={Traceability Matrix Showing the Connections Between Items and Other Sections}]
{colspec={l l l l l l l l l l l l l l l l l}, rowhead=1, hline{1,Z}=\heavyrulewidth, hline{2}=\lightrulewidth}
\textbf{} & \textbf{\hyperref[DD:qOverm]{DD:qOverm}} & \textbf{\hyperref[DD:initState]{DD:initState}} & \textbf{\hyperref[DD:eField]{DD:eField}} & \textbf{\hyperref[DD:bField]{DD:bField}} & \textbf{\hyperref[DD:regionRect]{DD:regionRect}} & \textbf{\hyperref[DD:fieldsByRegion]{DD:fieldsByRegion}} & \textbf{\hyperref[DD:detectorLine]{DD:detectorLine}} & \textbf{\hyperref[TM:velocity]{TM:velocity}} & \textbf{\hyperref[TM:acceleration]{TM:acceleration}} & \textbf{\hyperref[TM:lorentzForce]{TM:lorentzForce}} & \textbf{\hyperref[TM:eqnMotion]{TM:eqnMotion}} & \textbf{\hyperref[GD:kin2D]{GD:kin2D}} & \textbf{\hyperref[GD:vCrossB2D]{GD:vCrossB2D}} & \textbf{\hyperref[GD:dyn2D]{GD:dyn2D}} & \textbf{\hyperref[IM:stateEvol]{IM:stateEvol}} & \textbf{\hyperref[IM:detHit]{IM:detHit}}
\\
\hyperref[DD:qOverm]{DD:qOverm} &  &  &  &  &  &  &  &  &  &  &  &  &  &  &  & 
\\
\hyperref[DD:initState]{DD:initState} &  &  &  &  &  &  &  &  &  &  &  &  &  &  &  & 
\\
\hyperref[DD:eField]{DD:eField} &  &  &  &  &  &  &  &  &  &  &  &  &  &  &  & 
\\
\hyperref[DD:bField]{DD:bField} &  &  &  &  &  &  &  &  &  &  &  &  &  &  &  & 
\\
\hyperref[DD:regionRect]{DD:regionRect} &  &  &  &  &  &  &  &  &  &  &  &  &  &  &  & 
\\
\hyperref[DD:fieldsByRegion]{DD:fieldsByRegion} &  &  &  &  & X &  &  &  &  &  &  &  &  &  &  & 
\\
\hyperref[DD:detectorLine]{DD:detectorLine} &  &  &  &  &  &  &  &  &  &  &  &  &  &  &  & 
\\
\hyperref[TM:velocity]{TM:velocity} &  &  &  &  &  &  &  &  &  &  &  &  &  &  &  & 
\\
\hyperref[TM:acceleration]{TM:acceleration} &  &  &  &  &  &  &  &  &  &  &  &  &  &  &  & 
\\
\hyperref[TM:lorentzForce]{TM:lorentzForce} &  &  &  &  &  &  &  &  &  &  &  &  &  &  &  & 
\\
\hyperref[TM:eqnMotion]{TM:eqnMotion} &  &  &  &  &  &  &  &  &  & X &  &  &  &  &  & 
\\
\hyperref[GD:kin2D]{GD:kin2D} &  &  &  &  &  &  &  & X &  &  &  &  &  &  &  & 
\\
\hyperref[GD:vCrossB2D]{GD:vCrossB2D} &  &  &  &  &  &  &  &  &  &  &  &  &  &  &  & 
\\
\hyperref[GD:dyn2D]{GD:dyn2D} & X &  &  &  &  &  &  &  &  & X & X &  & X &  &  & 
\\
\hyperref[IM:stateEvol]{IM:stateEvol} & X & X &  &  &  & X &  &  &  &  &  & X &  & X &  & 
\\
\hyperref[IM:detHit]{IM:detHit} &  &  &  &  &  &  & X &  &  &  &  &  &  &  & X & 
\label{Table:TraceMatRefvsRef}
\end{longtblr}
\begin{longtblr}
[caption={Traceability Matrix Showing the Connections Between Requirements, Goal Statements and Other Items}]
{colspec={l l l l l l l l l l l l l l l l l l l l l l l l l l}, rowhead=1, hline{1,Z}=\heavyrulewidth, hline{2}=\lightrulewidth}
\textbf{} & \textbf{\hyperref[DD:qOverm]{DD:qOverm}} & \textbf{\hyperref[DD:initState]{DD:initState}} & \textbf{\hyperref[DD:eField]{DD:eField}} & \textbf{\hyperref[DD:bField]{DD:bField}} & \textbf{\hyperref[DD:regionRect]{DD:regionRect}} & \textbf{\hyperref[DD:fieldsByRegion]{DD:fieldsByRegion}} & \textbf{\hyperref[DD:detectorLine]{DD:detectorLine}} & \textbf{\hyperref[TM:velocity]{TM:velocity}} & \textbf{\hyperref[TM:acceleration]{TM:acceleration}} & \textbf{\hyperref[TM:lorentzForce]{TM:lorentzForce}} & \textbf{\hyperref[TM:eqnMotion]{TM:eqnMotion}} & \textbf{\hyperref[GD:kin2D]{GD:kin2D}} & \textbf{\hyperref[GD:vCrossB2D]{GD:vCrossB2D}} & \textbf{\hyperref[GD:dyn2D]{GD:dyn2D}} & \textbf{\hyperref[IM:stateEvol]{IM:stateEvol}} & \textbf{\hyperref[IM:detHit]{IM:detHit}} & \textbf{\hyperref[inputValues]{FR:Input-Values}} & \textbf{\hyperref[echoInputs]{FR:Echo-Inputs}} & \textbf{\hyperref[buildRegionGrid]{FR:Build-Region-Grid}} & \textbf{\hyperref[lookupActiveRegion]{FR:Lookup-Active-Region}} & \textbf{\hyperref[computeTraj]{FR:Compute-Trajectory}} & \textbf{\hyperref[reportOutputs]{FR:Report-Outputs}} & \textbf{\hyperref[correct]{NFR:Correctness}} & \textbf{\hyperref[maintainable]{NFR:Maintainability}} & \textbf{\hyperref[portable]{NFR:Portability}}
\\
\hyperref[predictTrajectory]{GS:predictTrajectory} &  &  &  &  &  &  &  &  &  &  &  &  &  &  & X &  &  &  &  &  &  &  &  &  & 
\\
\hyperref[determineDetectorOutcome]{GS:determineDetectorOutcome} &  &  &  &  &  &  &  &  &  &  &  &  &  &  &  & X &  &  &  &  &  &  &  &  & 
\\
\hyperref[inputValues]{FR:Input-Values} &  &  &  &  &  &  &  &  &  &  &  &  &  &  &  &  &  &  &  &  &  &  &  &  & 
\\
\hyperref[echoInputs]{FR:Echo-Inputs} &  &  &  &  &  &  &  &  &  &  &  &  &  &  &  &  &  &  &  &  &  &  &  &  & 
\\
\hyperref[buildRegionGrid]{FR:Build-Region-Grid} &  &  &  &  & X &  &  &  &  &  &  &  &  &  &  &  &  &  &  &  &  &  &  &  & 
\\
\hyperref[lookupActiveRegion]{FR:Lookup-Active-Region} &  &  &  &  &  & X &  &  &  &  &  &  &  &  &  &  &  &  &  &  &  &  &  &  & 
\\
\hyperref[computeTraj]{FR:Compute-Trajectory} &  &  &  &  &  &  &  &  &  &  &  &  &  &  & X &  &  &  &  &  &  &  &  &  & 
\\
\hyperref[reportOutputs]{FR:Report-Outputs} &  &  &  &  &  &  &  &  &  &  &  &  &  &  & X & X &  &  &  &  &  &  &  &  & 
\\
\hyperref[correct]{NFR:Correctness} &  &  &  &  &  &  &  &  &  &  &  &  &  &  &  &  &  &  &  &  &  &  &  &  & 
\\
\hyperref[maintainable]{NFR:Maintainability} &  &  &  &  &  &  &  &  &  &  &  &  &  &  &  &  &  &  &  &  &  &  &  &  & 
\\
\hyperref[portable]{NFR:Portability} &  &  &  &  &  &  &  &  &  &  &  &  &  &  &  &  &  &  &  &  &  &  &  &  & 
\label{Table:TraceMatAllvsR}
\end{longtblr}
The purpose of the traceability graphs is also to provide easy references on what has to be additionally modified if a certain component is changed. The arrows in the graphs represent dependencies. The component at the tail of an arrow is depended on by the component at the head of that arrow. Therefore, if a component is changed, the components that it points to should also be changed. \hyperref[Figure:TraceGraphAvsA]{Fig:TraceGraphAvsA} shows the dependencies of assumptions on each other. \hyperref[Figure:TraceGraphAvsAll]{Fig:TraceGraphAvsAll} shows the dependencies of data definitions, theoretical models, general definitions, instance models, requirements, likely changes, and unlikely changes on the assumptions. \hyperref[Figure:TraceGraphRefvsRef]{Fig:TraceGraphRefvsRef} shows the dependencies of data definitions, theoretical models, general definitions, and instance models on each other. \hyperref[Figure:TraceGraphAllvsR]{Fig:TraceGraphAllvsR} shows the dependencies of requirements and goal statements on the data definitions, theoretical models, general definitions, and instance models. \hyperref[Figure:TraceGraphAllvsAll]{Fig:TraceGraphAllvsAll} shows the dependencies of dependencies of assumptions, models, definitions, requirements, goals, and changes with each other.

\begin{figure}[H]
\begin{center}
\includesvg[width=\textwidth, inkscapelatex = false]{../../../../traceygraphs/trajecto/avsa}
\caption{TraceGraphAvsA}
\label{Figure:TraceGraphAvsA}
\end{center}
\end{figure}
\begin{figure}[H]
\begin{center}
\includesvg[width=\textwidth, inkscapelatex = false]{../../../../traceygraphs/trajecto/avsall}
\caption{TraceGraphAvsAll}
\label{Figure:TraceGraphAvsAll}
\end{center}
\end{figure}
\begin{figure}[H]
\begin{center}
\includesvg[width=\textwidth, inkscapelatex = false]{../../../../traceygraphs/trajecto/refvsref}
\caption{TraceGraphRefvsRef}
\label{Figure:TraceGraphRefvsRef}
\end{center}
\end{figure}
\begin{figure}[H]
\begin{center}
\includesvg[width=\textwidth, inkscapelatex = false]{../../../../traceygraphs/trajecto/allvsr}
\caption{TraceGraphAllvsR}
\label{Figure:TraceGraphAllvsR}
\end{center}
\end{figure}
\begin{figure}[H]
\begin{center}
\includesvg[width=\textwidth, inkscapelatex = false]{../../../../traceygraphs/trajecto/allvsall}
\caption{TraceGraphAllvsAll}
\label{Figure:TraceGraphAllvsAll}
\end{center}
\end{figure}
For convenience, the following graphs can be found at the links below:

\begin{itemize}
\item{\hyperref{../../../../traceygraphs/trajecto/avsa.svg}{}{}{TraceGraphAvsA}}
\item{\hyperref{../../../../traceygraphs/trajecto/avsall.svg}{}{}{TraceGraphAvsAll}}
\item{\hyperref{../../../../traceygraphs/trajecto/refvsref.svg}{}{}{TraceGraphRefvsRef}}
\item{\hyperref{../../../../traceygraphs/trajecto/allvsr.svg}{}{}{TraceGraphAllvsR}}
\item{\hyperref{../../../../traceygraphs/trajecto/allvsall.svg}{}{}{TraceGraphAllvsAll}}
\end{itemize}
\section{Values of Auxiliary Constants}
\label{Sec:AuxConstants}
There are no auxiliary constants.

\section{References}
\label{Sec:References}
\begin{filecontents*}{bibfile.bib}
@mastersthesis{koothoor2013,
author={Koothoor, Nirmitha},
title={A Document Driven Approach to Certifying Scientific Computing Software},
school={McMaster University},
year={2013},
address={Hamilton, ON, Canada}}
@article{parnasClements1986,
author={Parnas, David L. and Clements, P. C.},
title={A rational design process: How and why to fake it},
journal={IEEE Transactions on Software Engineering},
year={1986},
month=feb,
volume={12},
number={2},
pages={251--257},
address={Washington, USA}}
@article{smithKoothoor2016,
author={Smith, W. Spencer and Koothoor, Nirmitha},
title={A Document-Driven Method for Certifying Scientific Computing Software for Use in Nuclear Safety Analysis},
journal={Nuclear Engineering and Technology},
year={2016},
month=apr,
volume={48},
number={2},
pages={404--418},
howpublished={\url{http://www.sciencedirect.com/science/article/pii/S1738573315002582}}}
@inproceedings{smithLai2005,
author={Smith, W. Spencer and Lai, Lei},
title={A new requirements template for scientific computing},
booktitle={Proceedings of the First International Workshop on Situational Requirements Engineering Processes - Methods, Techniques and Tools to Support Situation-Specific Requirements Engineering Processes, SREP'05},
year={2005},
editor={Agerfalk, PJ and Kraiem, N. and Ralyte, J.},
address={Paris, France},
pages={107--121},
note={In conjunction with 13th IEEE International Requirements Engineering Conference,}}
@article{smithEtAl2007,
author={Smith, W. Spencer and Lai, Lei and Khedri, Ridha},
title={Requirements Analysis for Engineering Computation: A Systematic Approach for Improving Software Reliability},
journal={Reliable Computing, Special Issue on Reliable Engineering Computation},
year={2007},
month=feb,
volume={13},
number={1},
pages={83--107},
howpublished={\url{https://doi.org/10.1007/s11155-006-9020-7}}}
@misc{accelerationWiki,
author={Wikipedia Contributors},
title={Acceleration},
howpublished={\url{https://en.wikipedia.org/wiki/Acceleration}},
month=jun,
year={2019}}
@misc{velocityWiki,
author={Wikipedia Contributors},
title={Velocity},
howpublished={\url{https://en.wikipedia.org/wiki/Velocity}},
month=jun,
year={2019}}
\end{filecontents*}
\nocite{*}
\bibstyle{ieeetr}
\printbibliography[heading=none]
\end{document}
